%---------------------------------------------------------------
% SLO: slovenski povzetek
% ENG: slovenian abstract
%---------------------------------------------------------------
\selectlanguage{slovene} % Preklopi na slovenski jezik
\addcontentsline{toc}{chapter}{Povzetek}
\chapter*{Povzetek}

\noindent\textbf{Naslov:} \ttitle
\bigskip

Modeli za samodejno označevanje faz spanja se zaradi razlik v 
elektroencefalografskih (EEG) signalih pogosto slabo posplošujejo med različne
demografske skupine. V tem magistrskem delu predlagamo in ovrednotimo štiri
ogrodja na podlagi transformerjev, ki uporabljajo različne strategije obravnave
starosti: osnovni model brez podatka o starosti, eksplicitno pogojevanje z
vdelavo starosti v klasifikacijski žeton, ločeno naučene starostno specifične modele
ter prenosno učenje iz modela, naučenega na odraslih,
na ostale starostne skupine. Naši modeli so dosegli najvišje točnosti 76,8\,\% 
pri odraslih, 83,2\,\% pri najstnikih, 76,2\,\% pri otrocih in 74,8\,\% pri dojenčkih.
Naši pristopi za obravnavo starosti so presegli uveljavljeno klinično orodje YASA,
zlasti pri mlajših starostnih skupinah.

\subsection*{Ključne besede}
\textit{\tkeywords}
\clearemptydoublepage